% AIIB23 dataset -- pulled in by sections/03-dataset.tex
% Citation key nan2024aiib23 lives in references.bib

\subsection{The AIIB23 dataset}  % second dataset
\label{sec:aiib23}  % lets you write \ref{sec:aiib23} elsewhere

The AIIB23 dataset (Airway-Informed Quantitative CT Imaging Biomarker for
Fibrotic Lung Disease 2023) was released as part of an official challenge held
in conjunction with MICCAI 2023 and hosted on the CodaLab platform
\cite{nan2024aiib23}. It is the first publicly available dataset dedicated to
airway segmentation in patients with fibrotic lung disease. Earlier public
airway datasets, such as EXACT'09, BAS and ATM'22, mainly contain scans from
healthy subjects or from patients with common pulmonary conditions such as
nodules or COVID-19.

Fibrotic lung disease is particularly challenging for automatic airway
segmentation for two reasons. First, honeycombing, the presence of small
cystic airspaces in the lung periphery, visually resembles distal airway
branches, so models often produce false-positive predictions known as airway
leakages. Second, traction bronchiectasis and the distortion of lung tissue
alter airway morphology, which frequently leads to missing or discontinuous
branches. AIIB23 was therefore designed to benchmark segmentation methods on
these difficult cases.

\paragraph{Data sources and composition.}  % run-in bold heading
The dataset gathers high-resolution computed tomography (HRCT) scans from
three cohorts, acquired with different vendors and scanning protocols. The
training and validation sets come from the Open Source Imaging Consortium
(OSIC), an open-access database of patients with interstitial lung disease
(ILD), idiopathic pulmonary fibrosis (IPF) and other respiratory diseases.
The test set combines 90 fibrosis cases from the Australian IPF Repository
(AIPFR) with 50 COVID-19 cases from the University Hospital of Parma. The
COVID-19 cases make it possible to assess generalization to a different
pathology. The patient populations are broadly comparable in age, with a
median age of 68 years for OSIC, 70 for AIPFR and 65 for Parma.
Table~\ref{tab:aiib23_splits} summarizes the data splits.  % ~ = no line break before number

\begin{table}[htbp]  % float: let LaTeX place it (here, top, bottom, page)
  \centering  % center the table horizontally
  \caption{Composition of the AIIB23 dataset \cite{nan2024aiib23}.}  % caption above table
  \label{tab:aiib23_splits}  % target for \ref
  \begin{tabular}{llcc}  % 2 left-aligned + 2 centered columns
    \toprule  % thick top rule (booktabs)
    Split & Source & Task I (segmentation) & Task II (mortality) \\
    \midrule  % thin rule under the header
    Training   & OSIC          & 120 & 95 \\
    Validation & OSIC          & 52  & 52 \\
    Test       & AIPFR + Parma & 140\textsuperscript{a} & 90 \\
    \bottomrule  % thick bottom rule
  \end{tabular}

  \smallskip  % small gap before the note
  {\footnotesize \textsuperscript{a}90 fibrosis (AIPFR) and 50 COVID-19 (Parma) cases.}  % table note
\end{table}

\paragraph{Image selection and preprocessing.}  % run-in bold heading
All scans were converted from DICOM to NIfTI format using SimpleITK. Only
scans with more than 120 slices and an in-plane resolution of at least
$512 \times 512$ pixels were retained. Despite this criterion, the number of
slices varies considerably between cases. The fibrosis test scans contain
significantly fewer slices than the training data, which results in lower
image quality in the coronal and sagittal planes. This represents a domain
shift between training and test data.

\paragraph{Annotation protocol.}  % run-in bold heading
The ground-truth airway masks were produced through a semi-automatic,
multi-stage process. Preliminary labels were first generated by a deep
learning model trained on the BAS dataset. Three junior radiologists then
revised these labels, and two senior radiologists refined them by correcting
breakages and mis-annotated voxels. When annotators disagreed, the most
experienced radiologist's annotation was retained. Annotation was performed in
ITK-SNAP using a drawing tablet. Inter-observer agreement on airway branches
was very high (Cohen's $\kappa = 0.937$).

\paragraph{Tasks and evaluation.}  % run-in bold heading
The challenge defines two tasks. Task~I is binary airway segmentation of the
3D HRCT volume. It is evaluated with overlap metrics (intersection over union,
IoU, and precision) and with topological metrics that are particularly
relevant for thin tubular structures:
\begin{itemize}  % bulleted list of the branch metrics
  \item the detected length ratio (DLR) is the fraction of the ground-truth
        branch length that the prediction recovers;
  \item the detected branch ratio (DBR) is the fraction of ground-truth
        branches that are correctly detected, a branch counting as detected
        when its IoU with the corresponding ground-truth branch is at least 0.8;
  \item the airway leakage ratio (ALR) quantifies false-positive predictions.
\end{itemize}
The challenge's overall accuracy score is the mean of IoU, precision, DLR and
DBR. For reference, the best-performing challenge submission reached an IoU of
0.905, a DLR of 0.937 and a DBR of 0.905 on the test set.

% ---- OPTIONAL: keep this paragraph only if you used Task II ----
Task~II is a binary classification problem: predicting whether a patient is
alive or deceased 63 weeks after the first HRCT scan. The training set
contains 59 survivors and 36 deceased patients, and the validation set is
balanced (26/26). The test set is strongly imbalanced, with about 87\,\% of
survivors. Performance is measured with accuracy, AUC, sensitivity,
specificity and F1-score.
% ---- END OPTIONAL ----
